% Estilo base (evita overflow)
\lstdefinestyle{codebase}{
  inputencoding=utf8,
  backgroundcolor=\color{codebg},
  frame=single, rulecolor=\color{codeborder}, framerule=0.6pt, framesep=6pt,
  numbers=left, numberstyle=\tiny, numbersep=8pt,
  xleftmargin=1.5em, xrightmargin=0em,
  basicstyle=\ttfamily\footnotesize,
  keywordstyle=\bfseries\color{codekw},
  commentstyle=\itshape\color{codecm},
  stringstyle=\color{codestr},
  identifierstyle=\color{codeid},
  showstringspaces=false, upquote=true, tabsize=2, keepspaces=true,
  breaklines=true,            % <--- quiebra líneas largas
  breakatwhitespace=false,    % permite cortar dentro de tokens si es necesario
  columns=fullflexible,       % mejora el ajuste para evitar desbordes
  postbreak=\mbox{\textcolor{codecm}{\(\hookrightarrow\)}\space}, % marca de continuación
  captionpos=b, aboveskip=1ex, belowskip=1ex
}

% Estilo específico para C
\lstdefinestyle{cstyle}{
  style=codebase,
  language=C,
  morekeywords={uint8_t,uint16_t,uint32_t,bool,true,false}
}

% MATLAB/Octave (listings suele traer 'Matlab'; si no, usa 'Octave')
\lstdefinestyle{matlabstyle}{
  style=codebase,
  language=Matlab        % o: language=Octave
}

%  Mapea caracteres españoles en comentarios si hiciera falta:
\lstset{
  literate=
   {á}{{\'a}}1 {é}{{\'e}}1 {í}{{\'i}}1 {ó}{{\'o}}1 {ú}{{\'u}}1
   {Á}{{\'A}}1 {É}{{\'E}}1 {Í}{{\'I}}1 {Ó}{{\'O}}1 {Ú}{{\'U}}1
   {ñ}{{\~n}}1 {Ñ}{{\~N}}1 {°}{{$^\circ$}}1
}